% =====================================================================
%  Цель и задание лабораторной работы.
% =====================================================================

\textbf{Цель:} изучить механизмы контроля доступа, предоставляемые
серверными приложениями и системами управления базами данных, и получить
практические навыки настройки политики безопасности (администратор --
пользователь -- гость) для СУБД и веб-приложений.

\vspace{1em}

\textbf{Задание}:
\begin{enumerate}
  \item Создать программу, которая позволяет установить веб-приложение
  или систему управления базами данных (СУБД).
  \item Создать программу, которая позволяет настроить механизм контроля
  доступа для выбранного веб-приложения или СУБД.
  \item Создать программу для проверки корректности работы настроенной
  политики безопасности выбранной СУБД или веб-приложения.
  \item Создать программу, которая позволяет удалить веб-приложение или
  СУБД, а также созданные ими объекты файловой системы.
  \item Оформить отчёт о проделанной работе.
\end{enumerate}

Политика безопасности должна быть настроена следующим образом: в системе
должны присутствовать администратор системы с полным доступом ко всем
подсистемам и максимальными привилегиями, пользователь с полными либо
частичными правами в подсистеме, предоставленной администратором, и гость
с правами только на чтение отдельных фрагментов, предоставленных
администратором.

Задание выполнено для документоориентированной СУБД \texttt{MongoDB}. В
виде отдельного консольного C\#-проекта (\texttt{Lab1.2-MongoDB})
реализован сценарий, последовательно выполняющий все четыре программы из
задания как этапы одного сценария: установка, настройка контроля
доступа, проверка и удаление. Все действия -- реальные команды Linux
(\texttt{docker}, \texttt{mongosh} и т.\,п.), запускаемые как внешние
процессы; .NET здесь выступает только оркестратором сценария, а не
исполнителем самой политики безопасности.
